% 源分章 PDF 第 19 页；第 12.6 节及致谢。
\section{讨论}

SATE 是一个值得关注、也可能尚未得到充分利用的因果参数。在 CRT 中，候选单位往往通过系统性选择来满足研究的纳入标准。随后，通常还会应用匹配算法，选出匹配最好的 $n/2$ 对（Balzer 等，2015）。因此，观测数据往往并非通过从某个假想的簇总体或配对簇总体中进行简单随机抽样而得到。在这种情形下，与 PATE 相比，SATE 仍是一个容易解释的量，而且无需对抽样机制作出进一步假设，就能被严格估计。虽然研究发现的可推广性，以及将其迁移到新情境，仍然具有重要政策意义，但 SATE 和 PATE 都不能直接回答这一问题；这些新情境很可能在重要方面不同于当前样本，也不同于该样本所抽取自的任何假想目标总体。较好的做法是，将可推广性和可迁移性作为独立研究问题处理，它们需要各自的识别结果与相应的最优估计量（Cole 和 Stuart，2010；Stuart 等，2011；Bareinboim 和 Pearl，2013；Schochet，2013；Hartman 等，2015）。最后，对样本效应进行估计时，所达到的精度和功效至少与条件效应或总体效应一样高。

\noindent\textbf{致谢。}本章报告的研究获得美国国立卫生研究院（NIH）国立过敏与传染病研究所（NIAID）艾滋病研究部资助，项目编号为 R01-AI074345、R37-AI051164、UM1AI069502 和 U01AI099959。本章内容仅由作者负责，不一定代表 NIH 的官方观点。
\clearpage
\endgroup
